This appendix identifies the location and contents of the implementation, partition manifests, and result artefacts used in the experiments. They are released so that every result reported in Chapter~\ref{chap:ResultsAndDiscussion} can be audited on the identical partitions, and so that other researchers can evaluate their own models on the two public benchmarks without the contamination present in the released protocols.

\section{Code and Result Artefacts}
\label{sec:codeAvailability}
The code used to conduct the experiments is publicly available in the repository \href{https://github.com/reeyan-afzal-khan/THESIS-MASTERS-UTP}{THESIS-MASTERS-UTP}. It is organised as a shared Python library (\texttt{dasfe}) and a numbered notebook series that orchestrates it, from environment setup and dataset inventory through leakage assessment, partitioning, feature extraction, fusion, selection, benchmarking, cross-dataset generalisation with the Tomasov--Cao label mapping, the CNN and recent-architecture baselines, the spatial-confound ablation, repeated-partition and bootstrap statistics, protocol attribution, and the head-to-head comparison with recent architectures.

The frozen partition of each dataset is released as a manifest keyed by sample identifier, which records each window's presence before guard-band removal, after guard-band removal, and after decimation (\refss{sec:chap3Implementation}). Every experiment reads its split membership from this manifest rather than re-deriving it. The per-experiment predictions, full-precision metrics, experiment registry, the assumption files recorded for each reimplemented baseline, and a deterministic script that regenerates the reported tables are stored alongside the code.

The stored results are those reported in Chapter~\ref{chap:ResultsAndDiscussion}. The released library also contains the exact normalisation and the invariance audit of \refss{sec:chap5LimResidue}, so executing it yields the corrected figures of Table~\ref{tab:residue_effect}.
